\section{What Existing Agent-Security Work Does Not Settle}

The surrounding field is already crowded in productive ways. Prompt-injection work shows that LLM-integrated applications can confuse untrusted data with instructions, especially when external content is retrieved or processed inside a larger application \parencite{greshake2023not}. Instruction-hierarchy work treats the same vulnerability as a priority problem: systems must learn to follow privileged instructions over lower-priority or untrusted strings \parencite{wallace2024instruction}. AgentDojo turns the issue into a dynamic evaluation environment for tool-using agents, with tasks, attacks, defences, and security test cases that expose how untrusted content can interfere with tool-mediated action \parencite{debenedetti2024agentdojo}.

Those contributions are not rivals to delegation assurance. They supply pieces of the evidence base. A threat taxonomy can tell us which attacks to consider; an instruction hierarchy can tell us which sources should dominate; an agent benchmark can tell us whether the system resists particular attacks. Delegation assurance asks a different question: what evidence would justify the broader claim that an action remained within delegated authority?

The same distinction applies to evaluation work. LLM-as-judge methods can reduce evaluation costs and support scalable grading, but they also create a second-order assurance problem: the judge's label can license a safety claim even when the judge has collapsed task success, instruction following, policy compliance, risk, refusal, and failure attribution \parencite{zheng2023judging,liu2023geval,zeng2024llmbar}. Instruction-following benchmarks can make some compliance judgments more reproducible, but they do not by themselves settle whether an external action was authorized or whether an audit trace is sufficient for reconstruction \parencite{zhou2023ifeval}.

The gap is not that existing work ignores security. The gap is that technical controls, benchmarks, and audit records are often evaluated separately. Delegation assurance treats them as parts of one evidentiary claim about justified machine action.
